\documentclass[Main.tex]{subfiles}
\begin{document}
\section{Một số lương thực, thực phẩm}
\subsection{Lý thuyết trọng tâm}
\begin{kienthuccannho}
	\subsubsection{Vai trò của lương thực, thực phẩm}
	\begin{itemize}
		\item Một số lương thực, thực phẩm thông dụng:
		\begin{itemize}
			\item \textbf{\textit{Lương thực}} như gạo, ngô, khoai, sắn, … có chứa tinh bột.
			\item \textbf{\textit{Thực phẩm}} như thịt, cá, trứng, sữa, … được dùng để làm các món ăn.
		\end{itemize}
		\item Thức ăn của con người ở dạng lương thực (ngũ cốc: lúa gạo, ngô, khoai, sắn, lúa mì) và thực phẩm (như thịt, cá, rau, củ, quả, …). Thức ăn được cơ thể chuyển hóa thành năng lượng và các chất dinh dưỡng cần thiết cho cơ thể.
		\item Lương thực -- thực phẩm dễ bị hỏng trong không khí do nấm và vi khuẩn phân hủy nếu không được bảo quản hoặc bảo quản không đúng cách.
		\item Cách bảo quản lương thực -- thực phẩm: đông lạnh, hút chân không, hun khói, sấy khô, …
	\end{itemize}
	\subsubsection{Các nhóm chất dinh dưỡng trong lương thực – thực phẩm}
	\begin{center}
		\renewcommand{\arraystretch}{1.15}
		\begin{tabular}{|C{0.2\linewidth}|L{0.65\linewidth}|}
			\hline
			\textbf{Chất dinh dưỡng} & \textbf{Vai trò} \\
			\hline
			Carbohydrate (tinh bột, đường)
			&
			\begin{itemize}
				\item Chứa tinh bột, đường, chất xơ, nguồn gốc chủ yếu từ thực vật.
				\item Nguồn cung cấp năng lượng chính cần thiết cho các hoạt động của cơ thể.
				\item Carbohydrate có nhiều trong lúa, ngô, khoai, sắn, …
			\end{itemize}
			\\
			\hline
			Protein (đạm)
			&
			\begin{itemize}
				\item Có vai trò cấu tạo, duy trì và phát triển cơ thể.
				\item Protein liên quan đến mọi chức năng sống của cơ thể và cần thiết cho sự chuyển hóa các chất dinh dưỡng.
				\item Protein có nhiều trong cá, thịt, trứng, sữa và các loại hạt đậu, đỗ, …
			\end{itemize}
			\\
			\hline
			Lipid (chất béo)
			&
			\begin{itemize}
				\item Nguồn dự trữ năng lượng trong cơ thể và có tác dụng chống lạnh.
				\item Lipid có nhiều trong sữa, thịt, cá, lạc, vừng, bơ, dầu thực vật, …
			\end{itemize}
			\\
			\hline
			Khoáng chất và vitamin
			&
			\begin{itemize}
				\item Chất khoáng gồm calcium, phosphorus, sắt (iron), kẽm (zinc), …
				\item Vitamin gồm A, B1, B2, C, D, E, …
				\item Chất khoáng và vitamin có vai trò nâng cao hệ miễn dịch, giúp cơ thể khỏe mạnh, phòng chống bệnh tật.
				\item Nguồn giàu khoáng chất và vitamin như hải sản, rau xanh, củ quả tươi, …
			\end{itemize}
			\\
			\hline
		\end{tabular}
	\end{center}
	\subsubsection{Sức khỏe và chế độ dinh dưỡng}
	\begin{itemize}
		\item Các loại thức ăn khác nhau cung cấp năng lượng và các chất dinh dưỡng khác nhau. Nhu cầu về năng lượng và chất dinh dưỡng của mỗi người là khác nhau, phụ thuộc vào lứa tuổi, giới tính, cân nặng và chiều cao, …
		\item Nếu ăn quá nhiều mà ít hoạt động thì thức ăn sẽ được dự trữ dưới dạng chất béo; nếu ăn ít, không đủ chất thì cơ thể sẽ bị suy dinh dưỡng.
		\item Một số chất cần cho cơ thể với lượng nhỏ như khoáng chất, vitamin nhưng rất quan trọng.
	\end{itemize}
\end{kienthuccannho}
\subsection{Bài tập}

%%%=================Phần bài tập tự luận===================%%%
\phan{Bài tập tự luận}
\textit{\large Thí sinh trả lời từ câu 1 đến câu 14}
\Opensolutionfile{ansbth}[Ans/LGBT-C03_B09_MOT_SO_LUONG_THUC_THUC_PHAM]
\Opensolutionfile{ansbt}[Ans/AnsBT-C03_B09_MOT_SO_LUONG_THUC_THUC_PHAM]

%%%=============BT_1=============%%%
\begin{bt}%[6K3H4-1]
	Cho các lương thực -- thực phẩm sau:
	\begin{center}
		\renewcommand{\arraystretch}{1}
		\begin{tabular}{C{0.14\linewidth}C{0.14\linewidth}C{0.14\linewidth}C{0.14\linewidth}C{0.14\linewidth}C{0.14\linewidth}}
			\includegraphics[width=0.95\linewidth]{Images/DE_CUONG_HOA_6_HKI/lt_lua_gao.jpg}
			& \includegraphics[width=0.95\linewidth]{Images/DE_CUONG_HOA_6_HKI/lt_ngo.jpg}
			& \includegraphics[width=0.95\linewidth]{Images/DE_CUONG_HOA_6_HKI/lt_khoai_lang.jpg}
			& \includegraphics[width=0.95\linewidth]{Images/DE_CUONG_HOA_6_HKI/lt_mia.jpg}
			& \includegraphics[width=0.95\linewidth]{Images/DE_CUONG_HOA_6_HKI/lt_cac_loai_qua.jpg}
			& \includegraphics[width=0.95\linewidth]{Images/DE_CUONG_HOA_6_HKI/lt_mat_ong.jpg}
			\\
			Lúa, gạo & Ngô & Khoai lang & Mía & Các loại quả & Mật ong \\
			\includegraphics[width=0.95\linewidth]{Images/DE_CUONG_HOA_6_HKI/lt_ca.jpg}
			& \includegraphics[width=0.95\linewidth]{Images/DE_CUONG_HOA_6_HKI/lt_thit.jpg}
			& \includegraphics[width=0.95\linewidth]{Images/DE_CUONG_HOA_6_HKI/lt_trung.jpg}
			& \includegraphics[width=0.95\linewidth]{Images/DE_CUONG_HOA_6_HKI/lt_dau_do.jpg}
			& \includegraphics[width=0.95\linewidth]{Images/DE_CUONG_HOA_6_HKI/lt_dau_thuc_vat.jpg}
			& \includegraphics[width=0.95\linewidth]{Images/DE_CUONG_HOA_6_HKI/lt_bo.jpg}
			\\
			Cá & Thịt & Trứng & Đậu, đỗ & Dầu thực vật & Bơ \\
			\includegraphics[width=0.95\linewidth]{Images/DE_CUONG_HOA_6_HKI/lt_mo_lon.jpg}
			& \includegraphics[width=0.95\linewidth]{Images/DE_CUONG_HOA_6_HKI/lt_lac.jpg}
			& \includegraphics[width=0.95\linewidth]{Images/DE_CUONG_HOA_6_HKI/lt_vung.jpg}
			& \includegraphics[width=0.95\linewidth]{Images/DE_CUONG_HOA_6_HKI/lt_sua.png}
			& \includegraphics[width=0.95\linewidth]{Images/DE_CUONG_HOA_6_HKI/lt_rau_xanh.jpg}
			& \includegraphics[width=0.95\linewidth]{Images/DE_CUONG_HOA_6_HKI/lt_san.jpg}
			\\
			Mỡ lợn & Lạc & Vừng & Sữa & Rau xanh & Sắn \\
		\end{tabular}
	\end{center}
	\begin{enumerate}
		\item[(a)] Lương thực, thực phẩm nào có nguồn gốc từ thực vật, từ động vật?
		\item[(b)] Lương thực, thực phẩm nào có thể ăn sống, phải nấu chín?
	\end{enumerate}
	\loigiai{
		\begin{itemize}
			\item (a) Nguồn gốc từ thực vật: gạo, ngô, khoai lang, mía, các loại quả, đậu đỗ, dầu thực vật, bơ, lạc, vừng, rau xanh, sắn.

			Nguồn gốc từ động vật: mật ong, cá, thịt, trứng, mỡ lợn, sữa.
			\item (b) Có thể ăn sống: mía, các loại quả, mật ong, sữa, rau xanh.

			Phải nấu chín: gạo, ngô, khoai lang, cá, thịt, trứng, đậu -- đỗ, dầu thực vật, bơ, mỡ lợn, lạc, vừng, sắn.
		\end{itemize}
	}
\end{bt}

%%%=============BT_2=============%%%
\begin{bt}%[6K3N4-1]
	Cho các từ/cụm từ: \textit{lương thực, thực phẩm, bảo quản, tươi sống, chế biến}. Hãy chọn từ/cụm từ phù hợp điền vào chỗ trống để hoàn thành các phát biểu sau:
	\begin{enumerate}
		\item[(a)] Gạo, ngô, khoai, sắn là các loại ..(1).. chính ở Việt Nam.
		\item[(b)] Thịt, cá, tôm là các ..(2).. thường được dùng trong các bữa ăn hằng ngày. Chúng được ..(3).. để trở thành các món ăn.
		\item[(c)] Các thực phẩm ở dạng ..(4).. như thịt, cá cần được ..(5).. ở nhiệt độ thích hợp để bảo đảm an toàn cũng như tăng thời gian sử dụng.
		\item[(d)] Các ..(6).., ..(7).. cung cấp năng lượng và các chất thiết yếu giúp chúng ta có một cơ thể khỏe mạnh.
	\end{enumerate}
	\loigiai{
		\begin{itemize}
			\item (a) $(1)$ lương thực.
			\item (b) $(2)$ thực phẩm, $(3)$ chế biến.
			\item (c) $(4)$ tươi sống, $(5)$ bảo quản.
			\item (d) $(6)$ lương thực, $(7)$ thực phẩm.
		\end{itemize}
	}
\end{bt}

%%%=============BT_3=============%%%
\begin{bt}%[6K3V4-1]
	Việt Nam là quốc gia sản xuất và xuất khẩu gạo hàng đầu thế giới.
	\begin{enumerate}
		\item[(a)] Gạo là lương thực hay thực phẩm?
		\item[(b)] Kể tên hai khu vực sản xuất lúa gạo chính ở Việt Nam.
		\item[(c)] Tại sao phải thu hoạch lúa đúng thời vụ?
	\end{enumerate}
	\loigiai{
		\begin{itemize}
			\item (a) Gạo là lương thực, cung cấp lượng lớn tinh bột cho con người.
			\item (b) Hai khu vực sản xuất lúa gạo chính là Đồng bằng sông Cửu Long và Đồng bằng sông Hồng.
			\item (c) Cần phải thu hoạch lúa đúng thời vụ để:
			\begin{itemize}
				\item Đảm bảo hạt gạo có chất lượng tốt nhất.
				\item Tránh bị hao phí thu hoạch, vì nếu thu hoạch khi lúa chín quá thì hạt lúa rơi rụng ra đất rất nhiều.
				\item Chuẩn bị đất, kịp thời làm vụ khác.
			\end{itemize}
		\end{itemize}
	}
\end{bt}

%%%=============BT_4=============%%%
\begin{bt}%[6K3H4-1]
	Hãy nối thông tin hai cột cho phù hợp với nhau.
	\begin{center}
		\renewcommand{\arraystretch}{1.15}
		\begin{tabular}{|C{0.18\linewidth}|L{0.66\linewidth}|}
			\hline
			A. Chất béo & $(1)$ Chúng có vai trò như nhiên liệu của cơ thể. Sự tiêu hóa chuyển hóa chúng thành một loại đường đơn giản gọi là glucose, được đốt cháy để cung cấp năng lượng cho cơ thể hoạt động. \\
			\hline
			B. Carbohydrate & $(2)$ Nhờ dự trữ chúng dưới da mà các chú gấu có thể chống rét trong mùa đông lạnh giá. \\
			\hline
			C. Chất xơ & $(3)$ Chúng có trong nhiều bộ phận của cơ thể động vật và con người như tóc, cơ, máu, da, … \\
			\hline
			D. Protein & $(4)$ Con người chỉ cần một lượng nhỏ nhóm chất này nhưng có tác dụng lớn đến quá trình trao đổi chất. \\
			\hline
			E. Vitamin & $(5)$ Chúng không cung cấp dinh dưỡng nhưng cần cho quá trình tiêu hóa. \\
			\hline
		\end{tabular}
	\end{center}
	\loigiai{
		Kết quả nối: A -- $(2)$, B -- $(1)$, C -- $(5)$, D -- $(3)$, E -- $(4)$.
	}
\end{bt}

%%%=============BT_5=============%%%
\begin{bt}%[6K3V4-4]
	Lương thực -- thực phẩm tươi sống dễ bị hỏng, đặc biệt trong môi trường nóng ẩm. Hãy thu thập một số thông tin về lương thực -- thực phẩm phổ biến theo mẫu sau:
	\begin{center}
		\renewcommand{\arraystretch}{1.15}
		\begin{tabular}{|C{0.06\linewidth}|L{0.2\linewidth}|L{0.29\linewidth}|L{0.29\linewidth}|}
			\hline
			\textbf{STT} & \textbf{Lương thực -- thực phẩm} & \textbf{Dấu hiệu hư hỏng} & \textbf{Cách bảo quản} \\
			\hline
			1 & Gạo & Biến đổi màu sắc, có mốc xanh trên bề mặt. & Bảo quản trong chum, vại; đặt nơi khô thoáng, tránh ẩm. \\
			\hline
			2 & Thịt & & \\
			\hline
			3 & Trứng & & \\
			\hline
			4 & Cá & & \\
			\hline
			5 & Rau & & \\
			\hline
			6 & Trái cây & & \\
			\hline
		\end{tabular}
	\end{center}
	\loigiai{
		\begin{center}
			\renewcommand{\arraystretch}{1.15}
			\begin{tabular}{|C{0.06\linewidth}|L{0.2\linewidth}|L{0.29\linewidth}|L{0.29\linewidth}|}
				\hline
				\textbf{STT} & \textbf{Lương thực -- thực phẩm} & \textbf{Dấu hiệu hư hỏng} & \textbf{Cách bảo quản} \\
				\hline
				1 & Gạo & Biến đổi màu sắc, có mốc xanh trên bề mặt. & Bảo quản trong chum, vại; đặt nơi khô thoáng, tránh ẩm. \\
				\hline
				2 & Thịt & Biến đổi màu sắc, có mùi hôi, thịt mềm nhũn, chảy nước. & Bảo quản trong tủ lạnh. \\
				\hline
				3 & Trứng & Chảy nước, có mùi thối. & Bảo quản trong tủ lạnh hoặc nơi khô thoáng, tránh ánh nắng trực tiếp. \\
				\hline
				4 & Cá & Biến đổi màu sắc, có mùi hôi, thịt mềm nhũn, chảy nước. & Ướp đá hoặc ướp muối; bảo quản trong tủ lạnh. \\
				\hline
				5 & Rau & Lá bị biến đổi màu (vàng úa), có mùi. & Bảo quản trong tủ lạnh. \\
				\hline
				6 & Trái cây & Chảy nước, mềm nhũn, có mùi hôi. & Sấy khô, ngâm đường hoặc muối; bảo quản trong tủ lạnh. \\
				\hline
			\end{tabular}
		\end{center}
	}
\end{bt}

%%%=============BT_6=============%%%
\begin{bt}%[6K3V4-6]
	Khẩu phần ăn có ảnh hưởng rất lớn tới sức khỏe và sự phát triển của cơ thể con người. Hãy cho biết:
	\begin{enumerate}
		\item[(a)] Khẩu phần ăn đầy đủ phải bao gồm những chất dinh dưỡng nào.
		\item[(b)] Để xây dựng khẩu phần ăn hợp lí, ta cần dựa vào những căn cứ nào.
	\end{enumerate}
	\loigiai{
		\begin{itemize}
			\item (a) Khẩu phần ăn đầy đủ phải có đủ các chất dinh dưỡng: protein, lipid, carbohydrate, vitamin và chất khoáng.
			\item (b) Khẩu phần ăn hợp lí là khẩu phần ăn:
			\begin{itemize}
				\item Đảm bảo đủ lượng thức ăn phù hợp với từng đối tượng.
				\item Đảm bảo đủ các thành phần dinh dưỡng hữu cơ, vitamin, muối khoáng.
				\item Đảm bảo cung cấp đủ năng lượng cho nhu cầu của cơ thể.
			\end{itemize}
		\end{itemize}
	}
\end{bt}

%%%=============BT_7=============%%%
\begin{bt}%[6K3C4-6]
	Trong khẩu phần ăn của Dũng ($13$ tuổi) gồm có: $350$ g carbohydrate, $100$ g lipid, $200$ g protein và nhiều loại vitamin, muối khoáng khác. Em hãy cho biết khẩu phần của Dũng đã hợp lí chưa và giải thích rõ vì sao, biết:
	\begin{itemize}
		\item Hiệu suất hấp thụ của cơ thể đối với carbohydrate là $90\%$, lipid là $80\%$, protein là $60\%$.
		\item Nhu cầu dinh dưỡng của nam $13 - 15$ tuổi là khoảng $2\,500 - 2\,600$ kcal/ngày.
		\item $1$ g carbohydrate tạo ra $4{,}3$ kcal; $1$ g lipid tạo ra $9{,}3$ kcal; $1$ g protein tạo ra $4{,}1$ kcal.
	\end{itemize}
	\loigiai{
		\begin{itemize}
			\item Khối lượng carbohydrate hấp thụ: $350 \cdot 90\% = 315$ g.

			Năng lượng sinh ra từ $315$ g carbohydrate: $315 \cdot 4{,}3 = 1\,354{,}5$ kcal.
			\item Khối lượng lipid hấp thụ: $100 \cdot 80\% = 80$ g.

			Năng lượng sinh ra từ $80$ g lipid: $80 \cdot 9{,}3 = 744$ kcal.
			\item Khối lượng protein hấp thụ: $200 \cdot 60\% = 120$ g.

			Năng lượng sinh ra từ $120$ g protein: $120 \cdot 4{,}1 = 492$ kcal.
			\item Tổng năng lượng hấp thụ trong ngày: $1\,354{,}5 + 744 + 492 = 2\,590{,}5$ kcal.
		\end{itemize}
		Giá trị $2\,590{,}5$ kcal nằm trong khoảng $2\,500 - 2\,600$ kcal/ngày.
		\\
		\textbf{\textit{Kết luận:}} khẩu phần ăn của bạn Dũng là hợp lí vì đủ năng lượng và đủ các chất dinh dưỡng cần thiết cho cơ thể.
	}
\end{bt}

%%%=============BT_8=============%%%
\begin{bt}%[6K3H4-1]
	Những lương thực -- thực phẩm nào giàu chất bột đường, chất béo, chất đạm, vitamin và chất khoáng? Hãy kể tên những sản phẩm được chế biến từ các loại lương thực -- thực phẩm đó theo mẫu sau:
	\begin{center}
		\renewcommand{\arraystretch}{1.15}
		\begin{tabular}{|C{0.06\linewidth}|L{0.24\linewidth}|L{0.24\linewidth}|L{0.3\linewidth}|}
			\hline
			\textbf{STT} & \textbf{Các nhóm chất thiết yếu} & \textbf{Lương thực -- thực phẩm} & \textbf{Sản phẩm chế biến} \\
			\hline
			1 & Chất bột, đường & Gạo & Cơm, cháo, bánh, … \\
			\hline
			2 & Chất béo & & \\
			\hline
			3 & Chất đạm & & \\
			\hline
			4 & Vitamin và chất khoáng & & \\
			\hline
		\end{tabular}
	\end{center}
	\loigiai{
		\begin{center}
			\renewcommand{\arraystretch}{1.15}
			\begin{tabular}{|C{0.06\linewidth}|L{0.24\linewidth}|L{0.24\linewidth}|L{0.3\linewidth}|}
				\hline
				\textbf{STT} & \textbf{Các nhóm chất thiết yếu} & \textbf{Lương thực -- thực phẩm} & \textbf{Sản phẩm chế biến} \\
				\hline
				1 & Chất bột, đường & Gạo; khoai & Cơm, cháo, bánh, …; bánh, chè, khoai luộc \\
				\hline
				2 & Chất béo & Mỡ động vật (lợn, gà, …); đậu nành & Mỡ ăn; dầu ăn, sữa đậu nành \\
				\hline
				3 & Chất đạm & Thịt lợn; trứng & Thịt kho, thịt rang; trứng muối, trứng luộc \\
				\hline
				4 & Vitamin và chất khoáng & Rau; trái cây & Rau luộc, rau xào; trái cây khô, mứt \\
				\hline
			\end{tabular}
		\end{center}
	}
\end{bt}

%%%=============BT_9=============%%%
\begin{bt}%[6K3V4-8]
	Lương thực -- thực phẩm được chế biến để sử dụng làm thức ăn.
	\begin{enumerate}
		\item[(a)] Ở gia đình em thường sử dụng các cách chế biến lương thực -- thực phẩm nào?
		\item[(b)] Kể một số việc cần làm khi chế biến lương thực -- thực phẩm để bảo đảm vệ sinh và an toàn thực phẩm.
	\end{enumerate}
	\loigiai{
		\begin{itemize}
			\item (a) Các cách chế biến thường được sử dụng: hấp, luộc, nướng, rang, rán/chiên, xào.
			\item (b) Một số việc cần làm để bảo đảm vệ sinh và an toàn thực phẩm:
			\begin{itemize}
				\item Sử dụng găng tay, khẩu trang, tạp dề.
				\item Sử dụng dao, thớt riêng khi chế biến thực phẩm sống và thực phẩm chín.
				\item Rửa tay sạch sau khi tiếp xúc thực phẩm sống.
				\item Nấu chín các thực phẩm trước khi ăn.
			\end{itemize}
		\end{itemize}
	}
\end{bt}

%%%=============BT_10=============%%%
\begin{bt}%[6K3N4-1]
	Điền từ in nghiêng dưới đây vào chỗ trống cho phù hợp:
	\textit{chất dinh dưỡng, chuyển hóa, thức ăn, năng lượng}.
	\\
	Mọi cơ thể sống đều cần chất dinh dưỡng. Thực vật sử dụng ánh sáng mặt trời, nước và khí carbon dioxide để cung cấp cho chúng năng lượng. Động vật phải lấy …$(1)$… thông qua ăn thức ăn. Hầu hết …$(2)$… của chúng là thực vật hoặc động vật khác. Sau khi ăn, thức ăn được tiêu hóa, xảy ra các quá trình …$(3)$… để biến thức ăn thành các chất cơ thể cần.
	\loigiai{
		\begin{itemize}
			\item $(1)$: chất dinh dưỡng.
			\item $(2)$: thức ăn.
			\item $(3)$: chuyển hóa.
		\end{itemize}
	}
\end{bt}

%%%=============BT_11=============%%%
\begin{bt}%[6K3V4-1]
	Ta đã biết, $100$ g ngô và $100$ g gạo đều sinh ra năng lượng là $1\,528$ kJ. Vậy tại sao ta không ăn ngô thay gạo?
	\loigiai{
		$100$ g gạo cung cấp năng lượng bằng $100$ g ngô nhưng người ta thường xuyên ăn gạo vì gạo dễ tiêu hóa hơn ngô. Ngoài ra, gạo còn chứa những dưỡng chất tốt cho cơ thể nhiều hơn so với ngô.
	}
\end{bt}

%%%=============BT_12=============%%%
\begin{bt}%[6K3V4-6]
	Em hãy ghi lại thực đơn ngày hôm qua của em và xếp các thức ăn đó theo nhóm chất (carbohydrate, protein, chất béo, chất khoáng, vitamin).
	\loigiai{
		Học sinh ghi lại các món ăn đã dùng và sắp xếp theo nhóm chất. Ví dụ:
		\begin{center}
			\renewcommand{\arraystretch}{1.15}
			\begin{tabular}{|L{0.24\linewidth}|L{0.18\linewidth}|L{0.2\linewidth}|L{0.24\linewidth}|}
				\hline
				\textbf{Nhóm chất} & \textbf{Sáng} & \textbf{Trưa} & \textbf{Tối} \\
				\hline
				Carbohydrate & Bánh mì & Cơm & Cơm \\
				\hline
				Protein & Trứng & Thịt kho & Cá rán \\
				\hline
				Chất béo & Sữa & Thịt mỡ & Dầu thực vật (để xào rau) \\
				\hline
				Vitamin và khoáng chất & Rau thơm & Rau xanh, hoa quả & Rau xanh, hoa quả \\
				\hline
			\end{tabular}
		\end{center}
	}
\end{bt}

%%%=============BT_13=============%%%
\begin{bt}%[6K3V4-6]
	Đọc đoạn thông tin sau và trả lời các yêu cầu dưới đây.
	\begin{center}
		\textbf{CHẾ ĐỘ ĂN UỐNG HỢP LÍ}
	\end{center}
	Ăn uống là một nhu cầu thiết yếu của cơ thể. Chế độ ăn uống hợp lí sẽ bảo đảm sự phát triển tốt của cơ thể, phòng tránh bệnh tật. Một số nguyên tắc về chế độ ăn uống hợp lí được các chuyên gia dinh dưỡng đưa ra, đó là:
	\begin{itemize}
		\item Ăn đa dạng nhiều loại (bảo đảm đủ bốn nhóm: chất bột đường, vitamin và chất khoáng, chất đạm, chất béo).
		\item Phối hợp thức ăn nguồn đạm động vật và đạm thực vật.
		\item Ăn phối hợp dầu thực vật và mỡ động vật hợp lí.
		\item Ăn rau quả hằng ngày.
	\end{itemize}
	\begin{enumerate}
		\item[(a)] Kể tên một số nguồn đạm động vật và đạm thực vật mà em biết.
		\item[(b)] Kể tên một số loại dầu thực vật và mỡ động vật mà em biết.
		\item[(c)] Kể tên một số loại rau quả được sử dụng trong các bữa ăn hằng ngày của em. Chúng được chế biến như thế nào để làm thực phẩm trong bữa ăn?
	\end{enumerate}
	\loigiai{
		\begin{itemize}
			\item (a) Một số nguồn đạm động vật: thịt bò, thịt gà, thịt lợn, …

			Một số nguồn đạm thực vật: đậu tương, đậu xanh, …
			\item (b) Một số loại dầu thực vật: dầu đậu tương, dầu vừng, dầu lạc, …

			Một số loại mỡ động vật: mỡ lợn (heo), mỡ gà, mỡ bò, …
			\item (c) Một số rau quả dùng hằng ngày: rau cải, rau muống, rau cần, quả cam, quả bưởi, quả táo, … Chúng được chế biến bằng cách luộc, xào, nấu canh hoặc ăn tươi sau khi rửa sạch, gọt vỏ.
		\end{itemize}
	}
\end{bt}

%%%=============BT_14=============%%%
\begin{bt}%[6K3V4-3]
	Hiện tượng ngộ độc thực phẩm tập thể ngày càng nhiều. Trong đó, có không ít vụ ngộ độc thực phẩm xảy ra trong trường học.
	\begin{enumerate}
		\item[(a)] Kể tên một vài vụ ngộ độc thực phẩm mà em biết.
		\item[(b)] Em hãy nêu một số nguyên nhân dẫn đến ngộ độc thực phẩm.
		\item[(c)] Khi bị ngộ độc thực phẩm em cần phải làm gì?
		\item[(d)] Làm thế nào để phòng ngừa ngộ độc thực phẩm?
	\end{enumerate}
	\loigiai{
		\begin{itemize}
			\item (a) Học sinh nêu được một số vụ ngộ độc thực phẩm đã xảy ra ở địa phương hoặc trên báo, đài.
			\item (b) Một số nguyên nhân gây ngộ độc thực phẩm:
			\begin{itemize}
				\item Thực phẩm quá hạn sử dụng.
				\item Thực phẩm nhiễm khuẩn.
				\item Thực phẩm nhiễm hóa chất độc hại.
				\item Thực phẩm được chế biến không đảm bảo quy trình vệ sinh.
			\end{itemize}
			\item (c) Khi bị ngộ độc thực phẩm cần:
			\begin{itemize}
				\item Dừng ăn ngay thực phẩm đó.
				\item Có thể kích thích họng để tạo phản ứng nôn, nôn ra hết thực phẩm đã dùng.
				\item Pha oresol với nước cho người bị ngộ độc uống để tránh mất nước.
				\item Nếu ngộ độc nặng cần phải đưa tới bệnh viện cấp cứu.
				\item Nên lưu lại mẫu thực phẩm để dễ tìm hiểu nguyên nhân độc tố khi cần.
			\end{itemize}
			\item (d) Để phòng ngừa ngộ độc thực phẩm cần lưu ý:
			\begin{itemize}
				\item Ăn thực phẩm có nguồn gốc rõ ràng, còn hạn sử dụng.
				\item Kiểm tra kĩ thực phẩm trước khi ăn.
				\item Bảo đảm thực phẩm đưa vào chế biến là thực phẩm sạch, không nhiễm hóa chất độc hại.
				\item Chế biến thực phẩm phải đảm bảo vệ sinh.
			\end{itemize}
		\end{itemize}
	}
\end{bt}
\Closesolutionfile{ansbt}
\Closesolutionfile{ansbth}

%%%=================Phần trắc nghiệm nhiều lựa chọn===================%%%
\phan{Bài tập trắc nghiệm nhiều lựa chọn}
\textit{\large Thí sinh trả lời từ câu 1 đến câu 25. Mỗi câu thí sinh chỉ chọn một phương án}
\Opensolutionfile{ansex}[Ans/LGEX-C03_B09_MOT_SO_LUONG_THUC_THUC_PHAM]
\Opensolutionfile{ans}[Ans/Ans-C03_B09_MOT_SO_LUONG_THUC_THUC_PHAM]

%%%=============EX_1=============%%%
\begin{ex}%[6K3N4-1]
	Trường hợp nào sau đây là lương thực?
	\choice
	{\True Ngô}
	{Trứng}
	{Rau xanh}
	{Cá}
	\loigiai{Ngô là cây lương thực vì chứa nhiều tinh bột. Trứng, rau xanh, cá là thực phẩm.}
\end{ex}

%%%=============EX_2=============%%%
\begin{ex}%[6K3N4-1]
	Trường hợp nào sau đây là lương thực?
	\choice
	{Dưa hấu}
	{\True Khoai}
	{Sữa}
	{Cá}
	\loigiai{Khoai chứa nhiều tinh bột nên là lương thực. Dưa hấu, sữa, cá là thực phẩm.}
\end{ex}

%%%=============EX_3=============%%%
\begin{ex}%[6K3N4-1]
	Trường hợp nào sau đây là lương thực?
	\choice
	{Quả cam}
	{Mật ong}
	{Thịt}
	{\True Lúa}
	\loigiai{Lúa là cây lương thực vì hạt lúa (gạo) chứa nhiều tinh bột.}
\end{ex}

%%%=============EX_4=============%%%
\begin{ex}%[6K3N4-1]
	Trường hợp nào sau đây là thực phẩm?
	\choice
	{Khoai}
	{Sắn}
	{\True Thịt gà}
	{Lúa}
	\loigiai{Thịt gà là thực phẩm dùng để chế biến món ăn. Khoai, sắn, lúa là lương thực.}
\end{ex}

%%%=============EX_5=============%%%
\begin{ex}%[6K3N4-1]
	Trường hợp nào sau đây là thực phẩm?
	\choice
	{\True Quả nho}
	{Sắn}
	{Ngô}
	{Khoai}
	\loigiai{Quả nho là thực phẩm. Sắn, ngô, khoai chứa nhiều tinh bột nên là lương thực.}
\end{ex}

%%%=============EX_6=============%%%
\begin{ex}%[6K3H4-1]
	Cây trồng nào sau đây không được xem là cây lương thực?
	\choice
	{Lúa gạo}
	{Ngô}
	{\True Mía}
	{Lúa mì}
	\loigiai{Mía cung cấp đường chứ không phải nguồn tinh bột chính nên không được xem là cây lương thực.}
\end{ex}

%%%=============EX_7=============%%%
\begin{ex}%[6K3N4-1]
	Sản phẩm nào dưới đây chứa nhiều tinh bột?
	\choice
	{\True Gạo}
	{Trứng}
	{Rau xanh}
	{Dầu ăn}
	\loigiai{Gạo chứa nhiều tinh bột. Trứng giàu protein, rau xanh giàu vitamin, dầu ăn giàu lipid.}
\end{ex}

%%%=============EX_8=============%%%
\begin{ex}%[6K3N4-1]
	Gạo sẽ cung cấp chất dinh dưỡng nào nhiều nhất cho cơ thể?
	\choice
	{\True Carbohydrate (chất đường, bột)}
	{Protein (chất đạm)}
	{Lipid (chất béo)}
	{Vitamin}
	\loigiai{Gạo chứa nhiều tinh bột nên cung cấp chủ yếu carbohydrate cho cơ thể.}
\end{ex}

%%%=============EX_9=============%%%
\begin{ex}%[6K3N4-1]
	Trong các thực phẩm dưới đây, loại nào chứa nhiều protein (chất đạm) nhất?
	\choice
	{Gạo}
	{Ngô}
	{\True Thịt}
	{Gạo và rau xanh}
	\loigiai{Protein có nhiều trong cá, thịt, trứng, sữa và các loại hạt đậu, đỗ.}
\end{ex}

%%%=============EX_10=============%%%
\begin{ex}%[6K3N4-1]
	Sản phẩm nào dưới đây chứa nhiều chất béo?
	\choice
	{Thịt}
	{Trứng}
	{Rau xanh}
	{\True Dầu ăn}
	\loigiai{Dầu ăn là chất béo (lipid) ở dạng gần như nguyên chất nên chứa nhiều chất béo nhất.}
\end{ex}

%%%=============EX_11=============%%%
\begin{ex}%[6K3N4-1]
	Sản phẩm nào dưới đây chứa nhiều vitamin?
	\choice
	{Gạo}
	{Trứng}
	{\True Hoa quả}
	{Ngô}
	\loigiai{Hoa quả tươi là nguồn giàu vitamin và khoáng chất.}
\end{ex}

%%%=============EX_12=============%%%
\begin{ex}%[6K3H4-1]
	Lứa tuổi từ $11$ đến $15$ là lứa tuổi có sự phát triển nhanh chóng về chiều cao. Chất quan trọng nhất cho sự phát triển của xương là
	\choice
	{carbohydrate}
	{protein}
	{\True calcium}
	{chất béo}
	\loigiai{Calcium là chất khoáng cấu tạo nên xương và răng nên rất quan trọng cho sự phát triển chiều cao.}
\end{ex}

%%%=============EX_13=============%%%
\begin{ex}%[6K3H4-1]
	Thành phần dinh dưỡng chính của lương thực -- thực phẩm là
	\choice
	{carbohydrate, protein, khoáng chất}
	{tinh bột, đường, chất béo}
	{tinh bột, chất béo, vitamin}
	{\True carbohydrate, protein, chất béo, khoáng chất và vitamin}
	\loigiai{Lương thực -- thực phẩm cung cấp bốn nhóm chất dinh dưỡng chính: carbohydrate, protein, lipid (chất béo), khoáng chất và vitamin.}
\end{ex}

%%%=============EX_14=============%%%
\begin{ex}%[6K3H4-4]
	Phát biểu nào dưới đây đúng?
	\choice
	{Trong thành phần của ngô, khoai, sắn không chứa tinh bột}
	{\True Bảo quản thực phẩm không đúng cách làm giảm chất lượng thực phẩm}
	{Thực phẩm bị biến đổi tính chất vẫn sử dụng được}
	{Các thực phẩm phải nấu chín mới sử dụng được}
	\loigiai{Ngô, khoai, sắn đều chứa tinh bột. Thực phẩm bị biến đổi tính chất thì không nên sử dụng. Một số thực phẩm như hoa quả, rau sống có thể ăn sống. Do đó chỉ có phát biểu về bảo quản là đúng.}
\end{ex}

%%%=============EX_15=============%%%
\begin{ex}%[6K3H4-1]
	Trong các nhóm chất sau, những nhóm chất nào cung cấp năng lượng cho cơ thể?
	\begin{enumerate}[(1)]
		\item Chất đạm.
		\item Chất béo.
		\item Tinh bột, đường.
		\item Chất khoáng.
	\end{enumerate}
	\choice
	{$(1)$, $(2)$ và $(4)$}
	{$(2)$, $(3)$ và $(4)$}
	{$(1)$, $(2)$, $(3)$ và $(4)$}
	{\True $(1)$, $(2)$ và $(3)$}
	\loigiai{Chất đạm, chất béo và tinh bột, đường đều cung cấp năng lượng. Chất khoáng không cung cấp năng lượng mà giúp nâng cao hệ miễn dịch, cấu tạo cơ thể.}
\end{ex}

%%%=============EX_16=============%%%
\begin{ex}%[6K3H4-4]
	Việc làm nào dưới đây không phải cách bảo quản lương thực -- thực phẩm đúng?
	\choice
	{Chế biến cá và để trong tủ lạnh}
	{\True Để thịt ngoài không khí trong thời gian dài}
	{Sấy khô các loại trái cây}
	{Ướp muối cho cá}
	\loigiai{Để thịt ngoài không khí lâu ngày sẽ bị nấm và vi khuẩn phân hủy, gây ôi thiu nên đây không phải cách bảo quản đúng.}
\end{ex}

%%%=============EX_17=============%%%
\begin{ex}%[6K3V4-1]
	Cho các vật thể: sắn, rau muống, nước cam, nước mắm, khoai, ngô, kem, lúa. Số vật thể là lương thực là
	\choice
	{$5$}
	{$6$}
	{$3$}
	{\True $4$}
	\loigiai{Lương thực gồm sắn, khoai, ngô, lúa vì đều chứa nhiều tinh bột. Vậy có $4$ vật thể là lương thực.}
\end{ex}

%%%=============EX_18=============%%%
\begin{ex}%[6K3V4-1]
	Cho các vật thể: cá, ngô, thịt, sữa, khoai, sắn, mía, củ cải, quả táo, mật ong. Số vật thể là thực phẩm là
	\choice
	{$5$}
	{$6$}
	{\True $7$}
	{$8$}
	\loigiai{Thực phẩm gồm cá, thịt, sữa, mía, củ cải, quả táo, mật ong. Ngô, khoai, sắn là lương thực. Vậy có $7$ vật thể là thực phẩm.}
\end{ex}

%%%=============EX_19=============%%%
\begin{ex}%[6K3N4-1]
	Lương thực là thức ăn chứa nhiều chất nào sau đây?
	\choice
	{Protein}
	{\True Tinh bột}
	{Chất béo}
	{Vitamin}
	\loigiai{Đặc điểm chung của lương thực như gạo, ngô, khoai, sắn là đều chứa nhiều tinh bột.}
\end{ex}

%%%=============EX_20=============%%%
\begin{ex}%[6K3N4-1]
	Nhóm chất dinh dưỡng nào là nguồn dự trữ năng lượng trong cơ thể và có tác dụng chống lạnh?
	\choice
	{Carbohydrate}
	{Protein}
	{\True Lipid}
	{Vitamin}
	\loigiai{Lipid (chất béo) là nguồn dự trữ năng lượng trong cơ thể và có tác dụng chống lạnh.}
\end{ex}

%%%=============EX_21=============%%%
\begin{ex}%[6K3N4-1]
	Chất dinh dưỡng nào có vai trò cấu tạo, duy trì và phát triển cơ thể?
	\choice
	{Carbohydrate}
	{\True Protein}
	{Lipid}
	{Chất xơ}
	\loigiai{Protein (đạm) có vai trò cấu tạo, duy trì và phát triển cơ thể, liên quan đến mọi chức năng sống.}
\end{ex}

%%%=============EX_22=============%%%
\begin{ex}%[6K3H4-4]
	Cách bảo quản nào sau đây thường được dùng cho cá tươi?
	\choice
	{Để ở nơi có ánh nắng trực tiếp}
	{\True Ướp muối hoặc ướp đá, bảo quản trong tủ lạnh}
	{Ngâm trong chậu nước ở nhiệt độ phòng}
	{Để ngoài không khí nhiều ngày}
	\loigiai{Cá tươi rất dễ hỏng nên được ướp muối, ướp đá hoặc bảo quản trong tủ lạnh để hạn chế nấm và vi khuẩn phát triển.}
\end{ex}

%%%=============EX_23=============%%%
\begin{ex}%[6K3H4-3]
	Nguyên nhân nào sau đây KHÔNG gây ngộ độc thực phẩm?
	\choice
	{Thực phẩm quá hạn sử dụng}
	{Thực phẩm nhiễm khuẩn}
	{\True Thực phẩm sạch, được nấu chín và bảo quản đúng cách}
	{Thực phẩm nhiễm hóa chất độc hại}
	\loigiai{Thực phẩm sạch, nấu chín và bảo quản đúng cách là điều kiện bảo đảm an toàn nên không gây ngộ độc.}
\end{ex}

%%%=============EX_24=============%%%
\begin{ex}%[6K3H4-4]
	Vì sao lương thực -- thực phẩm dễ bị hỏng khi để lâu trong không khí?
	\choice
	{Vì chúng bị bay hơi hết nước}
	{\True Vì bị nấm và vi khuẩn phân hủy}
	{Vì chúng tự chuyển thành chất khí}
	{Vì chúng hấp thụ khí oxygen và nở ra}
	\loigiai{Lương thực -- thực phẩm dễ bị hỏng trong không khí do nấm và vi khuẩn phân hủy nếu không được bảo quản đúng cách.}
\end{ex}

%%%=============EX_25=============%%%
\begin{ex}%[6K3V4-1]
	Cho các thực phẩm: mỡ lợn, lạc, vừng, rau cải, dầu thực vật, quả cam. Số thực phẩm chứa nhiều lipid (chất béo) là
	\choice
	{$2$}
	{$3$}
	{\True $4$}
	{$5$}
	\loigiai{Thực phẩm chứa nhiều lipid gồm mỡ lợn, lạc, vừng, dầu thực vật. Rau cải và quả cam giàu vitamin, khoáng chất. Vậy có $4$ thực phẩm.}
\end{ex}
\Closesolutionfile{ans}
\Closesolutionfile{ansex}
%\bangdapan{Ans-C03_B09_MOT_SO_LUONG_THUC_THUC_PHAM}

%%%=================Phần trắc nghiệm đúng sai===================%%%
\phan{Bài tập trắc nghiệm đúng sai}
\textit{\large Thí sinh trả lời từ câu 1 đến câu 10. Trong mỗi ý a), b), c), d) ở mỗi câu thí sinh chọn đúng hoặc sai}
\Opensolutionfile{ansex}[Ans/LGTF-C03_B09_MOT_SO_LUONG_THUC_THUC_PHAM]
\Opensolutionfile{ansbook}[Ansbook/AnsTF-C03_B09_MOT_SO_LUONG_THUC_THUC_PHAM]
\Opensolutionfile{ans}[Ans/Tempt-C03_B09_MOT_SO_LUONG_THUC_THUC_PHAM]
\setcounter{ex}{0}

%%%=============TF_1=============%%%
\begin{ex}%[6K3H4-1]
	Xét các phát biểu về lương thực -- thực phẩm.
	\choiceTF
	{Lương thực là thức ăn của con người có nguồn gốc từ thực vật và đều chứa vitamin}
	{Thực phẩm là thức ăn của con người có nguồn gốc từ động vật}
	{\True Lúa, ngô, khoai, sắn là lương thực của con người}
	{\True Cá, thịt, trứng, sữa là thực phẩm của con người}
	\loigiai{
		\begin{itemchoice}[F1,F2,T3,T4]
			\itemch Sai vì điểm chung của lương thực là đều chứa tinh bột, không phải vitamin.
			\itemch Sai vì có những thực phẩm có nguồn gốc từ thực vật như rau, hoa quả.
			\itemch Lúa, ngô, khoai, sắn đều chứa nhiều tinh bột nên là lương thực.
			\itemch Cá, thịt, trứng, sữa là các thực phẩm quen thuộc trong bữa ăn hằng ngày.
	\end{itemchoice}}
\end{ex}

%%%=============TF_2=============%%%
\begin{ex}%[6K3H4-1]
	Xét các phát biểu về carbohydrate (tinh bột, đường) và protein (đạm).
	\choiceTF
	{\True Carbohydrate là chất dinh dưỡng chứa tinh bột, đường, chất xơ, nguồn gốc chủ yếu từ thực vật}
	{Carbohydrate là nguồn cung cấp năng lượng chính cho cơ thể, có nhiều trong rau xanh và hoa quả}
	{\True Protein có vai trò cấu tạo, duy trì và phát triển cơ thể}
	{\True Protein liên quan đến mọi chức năng sống của cơ thể, có nhiều trong cá, thịt, trứng, sữa và các loại đậu, đỗ}
	\loigiai{
		\begin{itemchoice}[T1,F2,T3,T4]
			\itemch Carbohydrate chứa tinh bột, đường, chất xơ và có nguồn gốc chủ yếu từ thực vật.
			\itemch Sai vì carbohydrate có nhiều trong ngũ cốc như lúa, ngô, khoai, sắn, … chứ không phải rau xanh, hoa quả.
			\itemch Protein có vai trò cấu tạo, duy trì và phát triển cơ thể.
			\itemch Protein liên quan đến mọi chức năng sống, có nhiều trong cá, thịt, trứng, sữa, đậu đỗ.
	\end{itemchoice}}
\end{ex}

%%%=============TF_3=============%%%
\begin{ex}%[6K3H4-1]
	Xét các phát biểu về lipid (chất béo) và khoáng chất, vitamin.
	\choiceTF
	{\True Lipid là nguồn dự trữ năng lượng trong cơ thể và có tác dụng chống lạnh}
	{\True Lipid có nhiều trong sữa, thịt, cá, lạc, bơ, dầu thực vật}
	{Chất khoáng và vitamin là nguồn cung cấp năng lượng chính cho con người}
	{\True Các nguồn giàu khoáng chất và vitamin như hải sản, rau xanh, củ quả tươi}
	\loigiai{
		\begin{itemchoice}[T1,T2,F3,T4]
			\itemch Lipid là nguồn dự trữ năng lượng và có tác dụng chống lạnh.
			\itemch Lipid có nhiều trong sữa, thịt, cá, lạc, vừng, bơ, dầu thực vật.
			\itemch Sai vì chất khoáng và vitamin có vai trò nâng cao hệ miễn dịch; năng lượng chủ yếu do carbohydrate, protein và lipid cung cấp.
			\itemch Hải sản, rau xanh, củ quả tươi là nguồn giàu khoáng chất và vitamin.
	\end{itemchoice}}
\end{ex}

%%%=============TF_4=============%%%
\begin{ex}%[6K3H4-2]
	Xét các phát biểu về tính chất của lương thực -- thực phẩm.
	\choiceTF
	{\True Lương thực -- thực phẩm có thể ở dạng tươi sống (như rau củ, cá, tôm, …) hoặc đã qua chế biến (như cơm, cá rán, thức ăn đóng hộp, …)}
	{Lương thực -- thực phẩm khó bị hỏng khi để lâu ngoài không khí}
	{\True Để bảo quản lương thực -- thực phẩm được lâu, người ta dùng cách đông lạnh, hút chân không, hun khói, sấy khô}
	{Mỗi bữa ăn không nên ăn đồng thời cả lương thực và thực phẩm để tránh cơ thể hấp thụ quá nhiều chất}
	\loigiai{
		\begin{itemchoice}[T1,F2,T3,F4]
			\itemch Lương thực -- thực phẩm có thể ở dạng tươi sống hoặc đã qua chế biến.
			\itemch Sai vì lương thực -- thực phẩm dễ bị hỏng do nấm và vi khuẩn phân hủy.
			\itemch Đông lạnh, hút chân không, hun khói, sấy khô là các cách bảo quản phổ biến.
			\itemch Sai vì cần ăn đồng thời với lượng phù hợp để bổ sung đa dạng chất dinh dưỡng cho cơ thể.
	\end{itemchoice}}
\end{ex}

%%%=============TF_5=============%%%
\begin{ex}%[6K3H4-1]
	Xét các phát biểu về vai trò của lương thực, thực phẩm.
	\choiceTF
	{\True Thức ăn được cơ thể chuyển hóa thành năng lượng và các chất dinh dưỡng cần thiết}
	{\True Gạo, ngô, khoai, sắn là các loại lương thực chính ở Việt Nam}
	{Cơ thể người có thể tự tổng hợp đầy đủ chất dinh dưỡng mà không cần thức ăn}
	{\True Thịt, cá, trứng, sữa là các thực phẩm thường dùng trong bữa ăn hằng ngày}
	\loigiai{
		\begin{itemchoice}[T1,T2,F3,T4]
			\itemch Thức ăn được cơ thể chuyển hóa thành năng lượng và các chất dinh dưỡng cần thiết.
			\itemch Gạo, ngô, khoai, sắn là các loại lương thực chính ở Việt Nam.
			\itemch Sai vì con người phải lấy chất dinh dưỡng thông qua ăn thức ăn, không tự tổng hợp được.
			\itemch Thịt, cá, trứng, sữa là các thực phẩm thông dụng hằng ngày.
	\end{itemchoice}}
\end{ex}

%%%=============TF_6=============%%%
\begin{ex}%[6K3H4-4]
	Xét các phát biểu về bảo quản lương thực -- thực phẩm.
	\choiceTF
	{\True Đông lạnh, hút chân không, hun khói, sấy khô là các cách bảo quản lương thực -- thực phẩm}
	{\True Gạo nên được bảo quản trong chum, vại, đặt nơi khô thoáng, tránh ẩm}
	{Thịt tươi có thể để ngoài không khí nhiều ngày mà không bị hỏng}
	{Thực phẩm đã biến đổi màu sắc, có mùi hôi vẫn có thể sử dụng nếu nấu thật chín}
	\loigiai{
		\begin{itemchoice}[T1,T2,F3,F4]
			\itemch Đây là các cách bảo quản lương thực -- thực phẩm phổ biến.
			\itemch Gạo cần để nơi khô thoáng, tránh ẩm để không bị mốc.
			\itemch Sai vì thịt tươi để ngoài không khí sẽ nhanh bị vi khuẩn phân hủy, gây ôi thiu.
			\itemch Sai vì thực phẩm đã biến đổi tính chất có thể chứa độc tố, không nên sử dụng.
	\end{itemchoice}}
\end{ex}

%%%=============TF_7=============%%%
\begin{ex}%[6K3H4-1]
	Xét các phát biểu về khoáng chất và vitamin.
	\choiceTF
	{\True Chất khoáng gồm calcium, phosphorus, sắt (iron), kẽm (zinc), …}
	{\True Cơ thể chỉ cần một lượng nhỏ khoáng chất và vitamin nhưng chúng rất quan trọng}
	{\True Calcium là chất quan trọng cho sự phát triển của xương}
	{Vitamin là nguồn cung cấp năng lượng chính cho cơ thể}
	\loigiai{
		\begin{itemchoice}[T1,T2,T3,F4]
			\itemch Chất khoáng gồm calcium, phosphorus, sắt, kẽm, …
			\itemch Khoáng chất và vitamin cần với lượng nhỏ nhưng rất quan trọng với cơ thể.
			\itemch Calcium là thành phần cấu tạo xương, răng nên rất quan trọng cho sự phát triển của xương.
			\itemch Sai vì vitamin không cung cấp năng lượng; năng lượng chủ yếu từ carbohydrate, protein, lipid.
	\end{itemchoice}}
\end{ex}

%%%=============TF_8=============%%%
\begin{ex}%[6K3H4-6]
	Xét các phát biểu về chế độ dinh dưỡng và sức khỏe.
	\choiceTF
	{Chỉ cần ăn thật nhiều thịt là có được chế độ dinh dưỡng hợp lí}
	{\True Nhu cầu năng lượng của mỗi người phụ thuộc vào lứa tuổi, giới tính, cân nặng và chiều cao}
	{\True Ăn quá nhiều mà ít hoạt động thì phần thức ăn dư được dự trữ dưới dạng chất béo}
	{\True Ăn ít, không đủ chất trong thời gian dài có thể dẫn đến suy dinh dưỡng}
	\loigiai{
		\begin{itemchoice}[F1,T2,T3,T4]
			\itemch Sai vì chế độ dinh dưỡng hợp lí cần đủ và cân đối cả bốn nhóm chất dinh dưỡng.
			\itemch Nhu cầu năng lượng phụ thuộc vào lứa tuổi, giới tính, cân nặng và chiều cao.
			\itemch Ăn nhiều mà ít vận động thì thức ăn dư được dự trữ dưới dạng chất béo.
			\itemch Ăn thiếu chất kéo dài sẽ dẫn đến suy dinh dưỡng.
	\end{itemchoice}}
\end{ex}

%%%=============TF_9=============%%%
\begin{ex}%[6K3V4-3]
	Xét các phát biểu về an toàn vệ sinh thực phẩm.
	\choiceTF
	{Nên dùng chung dao, thớt cho thực phẩm sống và thực phẩm chín để tiết kiệm}
	{\True Cần rửa tay sạch sau khi tiếp xúc với thực phẩm sống}
	{\True Nên chọn thực phẩm có nguồn gốc rõ ràng, còn hạn sử dụng}
	{\True Khi bị ngộ độc thực phẩm cần dừng ngay việc ăn thực phẩm đó và đưa người bệnh tới cơ sở y tế nếu nặng}
	\loigiai{
		\begin{itemchoice}[F1,T2,T3,T4]
			\itemch Sai vì phải dùng dao, thớt riêng cho thực phẩm sống và chín để tránh lây nhiễm vi khuẩn.
			\itemch Rửa tay sạch sau khi tiếp xúc thực phẩm sống giúp phòng ngừa nhiễm khuẩn.
			\itemch Thực phẩm có nguồn gốc rõ ràng, còn hạn sử dụng giúp phòng ngừa ngộ độc.
			\itemch Khi ngộ độc cần dừng ăn ngay và đưa tới cơ sở y tế nếu tình trạng nặng.
	\end{itemchoice}}
\end{ex}

%%%=============TF_10=============%%%
\begin{ex}%[6K3H4-2]
	Xét các phát biểu về nguồn gốc của lương thực -- thực phẩm.
	\choiceTF
	{\True Mật ong, cá, thịt, trứng, sữa có nguồn gốc từ động vật}
	{\True Gạo, ngô, khoai lang, rau xanh có nguồn gốc từ thực vật}
	{Mía, các loại quả, mật ong đều phải nấu chín mới ăn được}
	{\True Sắn, khoai lang cần được nấu chín trước khi ăn}
	\loigiai{
		\begin{itemchoice}[T1,T2,F3,T4]
			\itemch Mật ong, cá, thịt, trứng, sữa đều có nguồn gốc từ động vật.
			\itemch Gạo, ngô, khoai lang, rau xanh đều có nguồn gốc từ thực vật.
			\itemch Sai vì mía, các loại quả và mật ong đều có thể ăn sống.
			\itemch Sắn, khoai lang cần nấu chín trước khi ăn để bảo đảm an toàn và dễ tiêu hóa.
	\end{itemchoice}}
\end{ex}
\Closesolutionfile{ans}
\Closesolutionfile{ansbook}
\Closesolutionfile{ansex}
%\bangdapanTF{AnsTF-C03_B09_MOT_SO_LUONG_THUC_THUC_PHAM}

%%%=================Phần trắc nghiệm trả lời ngắn===================%%%
\phan{Bài tập trắc nghiệm trả lời ngắn}
\textit{\large Thí sinh trả lời từ câu 1 đến câu 10}
\Opensolutionfile{ansex}[Ans/LGSA-C03_B09_MOT_SO_LUONG_THUC_THUC_PHAM]
\Opensolutionfile{ansexh}[Ans/AnsSA-C03_B09_MOT_SO_LUONG_THUC_THUC_PHAM]
\setcounter{ex}{0}

%%%=============SA_1=============%%%
\begin{ex}%[6K3N4-1]
	Cho các vật thể: $(1)$ lúa, $(2)$ mía, $(3)$ cá, $(4)$ rau xanh, $(5)$ sắn, $(6)$ khoai. Hãy sắp xếp các vật thể thuộc loại lương thực thành bộ số theo số thứ tự tăng dần (ví dụ $12$, $134$, …).
	\shortans{$156$}
	\loigiai{
		\begin{itemize}
			\item Lương thực gồm lúa $(1)$, sắn $(5)$, khoai $(6)$ vì đều chứa nhiều tinh bột.
			\item Mía, cá, rau xanh là thực phẩm.
		\end{itemize}
		\textbf{\textit{Đáp số:} $156$.}
	}
\end{ex}

%%%=============SA_2=============%%%
\begin{ex}%[6K3H4-2]
	Cho các vật thể: $(1)$ mỡ lợn, $(2)$ vừng, $(3)$ gạo, $(4)$ thịt, $(5)$ mật ong, $(6)$ quả nho. Hãy sắp xếp các lương thực -- thực phẩm cần phải nấu chín trước khi ăn thành bộ số theo số thứ tự tăng dần (ví dụ $12$, $134$, …).
	\shortans{$1234$}
	\loigiai{
		\begin{itemize}
			\item Cần nấu chín trước khi ăn: mỡ lợn $(1)$, vừng $(2)$, gạo $(3)$, thịt $(4)$.
			\item Mật ong và quả nho có thể ăn sống.
		\end{itemize}
		\textbf{\textit{Đáp số:} $1234$.}
	}
\end{ex}

%%%=============SA_3=============%%%
\begin{ex}%[6K3H4-1]
	Cho các phát biểu sau:
	\begin{enumerate}[(1)]
		\item Tinh bột, đường là nguồn cung cấp năng lượng chính cho các hoạt động của cơ thể.
		\item Protein có nhiều trong rau xanh, củ quả tươi.
		\item Lipid là nguồn dự trữ năng lượng trong cơ thể và có tác dụng chống lạnh.
		\item Cơ thể cần một lượng lớn khoáng chất và vitamin để duy trì hoạt động.
		\item Thực phẩm dễ bị ôi thiu nên cần bảo quản cẩn thận như cho vào tủ lạnh, hút chân không, …
	\end{enumerate}
	Hãy sắp xếp các phát biểu đúng thành bộ số theo số thứ tự tăng dần (ví dụ $12$, $134$, …).
	\shortans{$135$}
	\loigiai{
		\begin{itemize}
			\item $(2)$ sai vì protein có nhiều trong cá, thịt, trứng, sữa và các loại đậu, đỗ.
			\item $(4)$ sai vì cơ thể chỉ cần một lượng nhỏ khoáng chất và vitamin.
			\item Các phát biểu đúng là $(1)$, $(3)$, $(5)$.
		\end{itemize}
		\textbf{\textit{Đáp số:} $135$.}
	}
\end{ex}

%%%=============SA_4=============%%%
\begin{ex}%[6K3H4-1]
	Cho các thực phẩm: $(1)$ gạo, $(2)$ thịt bò, $(3)$ trứng, $(4)$ khoai lang, $(5)$ cá, $(6)$ ngô. Hãy sắp xếp các thực phẩm giàu protein (chất đạm) thành bộ số theo số thứ tự tăng dần (ví dụ $12$, $134$, …).
	\shortans{$235$}
	\loigiai{
		\begin{itemize}
			\item Giàu protein gồm thịt bò $(2)$, trứng $(3)$, cá $(5)$.
			\item Gạo, khoai lang, ngô giàu carbohydrate (tinh bột).
		\end{itemize}
		\textbf{\textit{Đáp số:} $235$.}
	}
\end{ex}

%%%=============SA_5=============%%%
\begin{ex}%[6K3H4-1]
	Cho các loại: $(1)$ mỡ lợn, $(2)$ gạo, $(3)$ dầu thực vật, $(4)$ lạc, $(5)$ rau cải, $(6)$ vừng. Hãy sắp xếp các loại giàu lipid (chất béo) thành bộ số theo số thứ tự tăng dần (ví dụ $12$, $134$, …).
	\shortans{$1346$}
	\loigiai{
		\begin{itemize}
			\item Giàu lipid gồm mỡ lợn $(1)$, dầu thực vật $(3)$, lạc $(4)$, vừng $(6)$.
			\item Gạo giàu tinh bột, rau cải giàu vitamin và chất khoáng.
		\end{itemize}
		\textbf{\textit{Đáp số:} $1346$.}
	}
\end{ex}

%%%=============SA_6=============%%%
\begin{ex}%[6K3N4-1]
	Cho các vật thể: $(1)$ lúa mì, $(2)$ trứng, $(3)$ khoai lang, $(4)$ sữa, $(5)$ ngô, $(6)$ tôm. Hãy sắp xếp các vật thể thuộc loại lương thực thành bộ số theo số thứ tự tăng dần (ví dụ $12$, $134$, …).
	\shortans{$135$}
	\loigiai{
		\begin{itemize}
			\item Lương thực gồm lúa mì $(1)$, khoai lang $(3)$, ngô $(5)$ vì đều chứa nhiều tinh bột.
			\item Trứng, sữa, tôm là thực phẩm.
		\end{itemize}
		\textbf{\textit{Đáp số:} $135$.}
	}
\end{ex}

%%%=============SA_7=============%%%
\begin{ex}%[6K3H4-2]
	Cho các lương thực -- thực phẩm: $(1)$ mía, $(2)$ gạo, $(3)$ mật ong, $(4)$ sắn, $(5)$ quả cam, $(6)$ trứng. Hãy sắp xếp các loại có thể ăn sống (không cần nấu chín) thành bộ số theo số thứ tự tăng dần (ví dụ $12$, $134$, …).
	\shortans{$135$}
	\loigiai{
		\begin{itemize}
			\item Có thể ăn sống gồm mía $(1)$, mật ong $(3)$, quả cam $(5)$.
			\item Gạo, sắn, trứng cần được nấu chín trước khi ăn.
		\end{itemize}
		\textbf{\textit{Đáp số:} $135$.}
	}
\end{ex}

%%%=============SA_8=============%%%
\begin{ex}%[6K3H4-4]
	Cho các cách xử lí sau: $(1)$ đông lạnh, $(2)$ để ngoài trời nắng nóng, $(3)$ hút chân không, $(4)$ hun khói, $(5)$ sấy khô. Hãy sắp xếp các cách bảo quản lương thực -- thực phẩm đúng thành bộ số theo số thứ tự tăng dần (ví dụ $12$, $134$, …).
	\shortans{$1345$}
	\loigiai{
		\begin{itemize}
			\item Các cách bảo quản đúng là đông lạnh $(1)$, hút chân không $(3)$, hun khói $(4)$, sấy khô $(5)$.
			\item $(2)$ sai vì để ngoài trời nắng nóng làm thực phẩm nhanh hỏng do nấm và vi khuẩn phát triển.
		\end{itemize}
		\textbf{\textit{Đáp số:} $1345$.}
	}
\end{ex}

%%%=============SA_9=============%%%
\begin{ex}%[6K3H4-1]
	Cho các phát biểu sau:
	\begin{enumerate}[(1)]
		\item Carbohydrate có nhiều trong lúa, ngô, khoai, sắn.
		\item Protein có nhiều trong cá, thịt, trứng, sữa, đậu đỗ.
		\item Lipid chỉ có trong thực vật, không có trong động vật.
		\item Khoáng chất và vitamin giúp nâng cao hệ miễn dịch của cơ thể.
		\item Lương thực -- thực phẩm không bị hỏng khi để lâu ngoài không khí.
	\end{enumerate}
	Hãy sắp xếp các phát biểu đúng thành bộ số theo số thứ tự tăng dần (ví dụ $12$, $134$, …).
	\shortans{$124$}
	\loigiai{
		\begin{itemize}
			\item $(3)$ sai vì lipid có cả trong động vật (mỡ lợn, bơ, sữa) và thực vật (lạc, vừng, dầu thực vật).
			\item $(5)$ sai vì lương thực -- thực phẩm dễ bị hỏng do nấm và vi khuẩn phân hủy.
			\item Các phát biểu đúng là $(1)$, $(2)$, $(4)$.
		\end{itemize}
		\textbf{\textit{Đáp số:} $124$.}
	}
\end{ex}

%%%=============SA_10=============%%%
\begin{ex}%[6K3H4-1]
	Cho các nguồn thực phẩm: $(1)$ hải sản, $(2)$ rau xanh, $(3)$ mỡ lợn, $(4)$ củ quả tươi, $(5)$ dầu ăn. Hãy sắp xếp các nguồn giàu khoáng chất và vitamin thành bộ số theo số thứ tự tăng dần (ví dụ $12$, $134$, …).
	\shortans{$124$}
	\loigiai{
		\begin{itemize}
			\item Nguồn giàu khoáng chất và vitamin gồm hải sản $(1)$, rau xanh $(2)$, củ quả tươi $(4)$.
			\item Mỡ lợn và dầu ăn là nguồn giàu lipid (chất béo).
		\end{itemize}
		\textbf{\textit{Đáp số:} $124$.}
	}
\end{ex}
\Closesolutionfile{ansexh}
\Closesolutionfile{ansex}
%\bangdapanSA{AnsSA-C03_B09_MOT_SO_LUONG_THUC_THUC_PHAM}
\end{document}
